%% ma: L12-B07-0031
%% lop: 12
%% chuong: 2
%% bai: 7
%% dang: 12.7.7
%% loai: TN4
%% muc: TH
%% dap_an: "A"
%% nguon: "CĐ2.6.pdf giáo viên gửi 27/10/2026 (M07, Lê Minh Khôi – Nguyễn Anh Khôi, Chương II Chủ đề 6 Ứng dụng vào các bài toán thực tế), Câu 11"
%% kien_thuc: "Bài 7 (tọa độ theo các hướng Đông, Tây, Nam, Bắc), Bài 8 (khoảng cách hai điểm)"
%% trang_thai: chinh-thuc
\begin{ex}
Hai chiếc khinh khí cầu cùng bay lên tại một địa điểm. Sau một thời gian bay, chiếc khinh khí cầu thứ nhất cách điểm xuất phát về phía Đông $100$ km và về phía Nam $80$ km, đồng thời cách mặt đất $1$ km. Chiếc khinh khí cầu thứ hai cách điểm xuất phát về phía Bắc $70$ km và về phía Tây $60$ km, đồng thời cách mặt đất $800$ m. Khoảng cách giữa hai khinh khí cầu là bao nhiêu?
\choice{\True $219{,}32$ km}{$210{,}00$ km}{$225{,}15$ km}{$205{,}45$ km}
\loigiai{Chọn $Ox$ hướng Nam, $Oy$ hướng Đông, $Oz$ hướng lên (đơn vị km): $K_1(80;100;1)$, $K_2(-70;-60;0{,}8)$.

$K_1K_2=\sqrt{150^2+160^2+0{,}2^2}\approx219{,}32$ km.}
\end{ex}
